% Generado por scripts/70_tablas_figuras.R a partir de resultados-v1. No editar a mano.
\begin{table}[htbp]
\centering
\small
\caption{Tabla 4.4—Características de línea base de los municipios tratados y de control}
\label{tab:tabla_4_4}
\begin{threeparttable}
\begin{tabular}{@{}lccc@{}}
\toprule
 & Tratados & Controles & Diferencia \\
\midrule
Pérdida de bosque, promedio 2001--2009 (ha/año) & 873,5 & 158,5 & 0,55 \\
 & (1.728,8) & (592,0) &  \\
Tendencia de la pérdida 2001--2009 (pendiente, ha por año) & 29,0 & 5,0 & 0,21 \\
 & (154,3) & (47,6) &  \\
Bosque en 2000 (miles de ha) & 497,9 & 51,0 & 0,77 \\
 & (801,0) & (154,7) &  \\
Temperatura media, 2001 (°C) & 22,9 & 19,5 & 0,79 \\
 & (3,5) & (5,0) &  \\
Precipitación anual, 2001 (mm) & 4.564 & 1.797 & 1,74 \\
 & (2.040) & (959) &  \\
Área en el SINAP público, 2012 (\%) & 8,9 & 6,5 & 0,15 \\
 & (16,1) & (15,9) &  \\
Distancia a la capital departamental & 103,7 & 80,0 & 0,36 \\
 & (70,0) & (59,5) &  \\
Distancia a Bogotá & 391,1 & 315,7 & 0,45 \\
 & (129,5) & (196,9) &  \\
Distancia al mercado mayorista & 168,7 & 63,7 & 0,82 \\
 & (154,5) & (96,1) &  \\
Cultivos de coca (ha) & 775,0 & 102,0 & 0,49 \\
 & (1.788,8) & (744,5) &  \\
Homicidios, 2003 & 18,6 & 16,2 & 0,04 \\
 & (48,3) & (77,4) &  \\
\midrule
Municipios & 36 & 986 &  \\
\bottomrule
\end{tabular}
\begin{tablenotes}[flushleft]
\footnotesize
\item \textit{Notas:} Medias por grupo; desviación estándar entre paréntesis. La diferencia es estandarizada: diferencia de medias dividida por la raíz del promedio de las varianzas. Muestra principal: municipios tratados de la D1 depurada con adopción entre 2013 y 2019 y controles limpios, sin ningún vínculo con proyectos de carbono. Las distancias y la tasa de homicidios están en las unidades del panel municipal del CEDE.
\item \textit{Fuente:} Global Forest Change v1.12; ERA5-Land y CHIRPS; RUNAP; panel municipal del CEDE.
\end{tablenotes}
\end{threeparttable}
\end{table}
